\chapter{HPD position reconstruction}
\label{chap:recon}

I need an introduction on the detector geometry before describing how we can make the position reconstruction.

We need to talk about the cluster size distribution and then we can talk about the implementation of different methods for various cluster sizes. We should also first point out that we are always excluding the fourth and successive pixels in an event, because with the given properties of the sensor, the charge collected by the fourth pixel is always below 20 e, while we currently have a higher noise level. This means that we cannot reconstruct the position using this information. 


The main feature of this detector is that the silicon sensor is on top of the readout chip, and its thickness is enough to convert all the energy of a soft X-ray photon inside it, while producing a small diffusion of the produced electrons, with respect to the size of the pixels. This can be seen with a simple calculation, because the transverse diffusion has a sigma of about 40 um/cm**0.5 in silicon, and with a sensor 300 um. This means that for this thickness, the charge cloud has a sigma of about 7 um, smaller than the pixel pitch.

The first method that can be thought to reconstruct the incident position, is to calculate the charge barycenter using the position of the hit pixels and the charge that they collected. It can be shown that this method is not accurate at all, showing a high degree of bias in the reconstruction. This is due to the charge diffusion nature and the electron cloud size. The cloud 98\% of the charge is contained in a 2 sigma circle, which means that a pixel can collect signal only if the photon hits at about this length from the pixel edge, otherwise no charge is collected. Charge barycentering is highly problematic when a photon hits in the edge region, where even if a few electron hit the second pixel, the position would be reconstructed in proximity of the first pixel center. This problem is evident in the case of two 